\documentclass[10pt]{article}
\usepackage[margin=0.9in]{geometry}
\usepackage{amsmath,amssymb,amsthm,microtype}
\usepackage[hidelinks]{hyperref}
\setlength{\emergencystretch}{3em}
\newtheorem{theorem}{Theorem}
\title{A Ternary Obstruction to Bounded-Odd-Step\\Inverse Collatz Searches}
\author{Collatz proof project}
\date{September 5, 2026}
\begin{document}
\maketitle
\begin{abstract}
For every fixed odd-step budget $K$, every positive-length shortcut Collatz
segment ending at a target congruent to one modulo $3^K$ and containing
at most $K$ odd steps has a contracting final coefficient. The total
segment length is unrestricted. Consequently no seed whose coefficients
never contract can reach that target within the budget. We prove this
in Lean using a sharp canonical-residue endpoint bound. Arbitrarily large
targets outside multiples of three have this obstruction, but none is
asserted to be a noncontracting seed. The result limits a proposed backward
exclusion method; it does not prove Collatz. No priority claim is made.
\end{abstract}

\section{The search question}
Let $T(n)=n/2$ for even $n$ and $(3n+1)/2$ for odd $n$. Let $j_t(n)$
count odd states before time $t$. A seed satisfies
$\operatorname{NeverContracts}(n)$ when $2^t\le3^{j_t(n)}$ for every $t$.
The preceding noncontracting-tail theorem proves that every unbounded
positive orbit has a tail in this class, and that every positive seed
in this class is unbounded.

A possible exclusion method is to take a least positive seed in the class
and construct a smaller predecessor that is also in the class. Backward
transport works through a prefix whose coefficients all stay at least one.
The present theorem shows why a fixed bound on the number of odd steps in
such constructions leaves a whole residue class unresolved, regardless
of how many even steps are permitted.

\section{A sharp endpoint bound}
\begin{theorem}\label{thm:canonical}
For every $t$ and every $0\le r<2^t$,
\[
T^t(r)<3^{j_t(r)}.
\]
\end{theorem}
\begin{proof}
Induct on $t$ using binary residue lifts. At depth $t$, write $e=T^t(r)$
and $A=3^{j_t(r)}$, with $e<A$. Adding the next binary digit replaces
$r$ by $r+2^t\varepsilon$, where $\varepsilon\in\{0,1\}$. The exact
affine transport identity gives the intermediate value
$p=e+A\varepsilon<2A$, with unchanged first-$t$ odd count.
If $p$ is even, its next value $p/2$ is below $A$. If $p$ is odd, then
$p\le2A-1$ and $(3p+1)/2\le3A-1<3A$; the odd count increases by one.
This proves the next-depth bound. At depth zero the only residue is zero.
\end{proof}

\section{The residue-one obstruction}
\begin{theorem}\label{thm:full}
If $t>0$ and $2^t\le3^{j_t(n)}$, then
\[
T^t(n)\not\equiv1\pmod{3^{j_t(n)}}.
\]
\end{theorem}
\begin{proof}
Put $r=n\bmod2^t$ and $A=3^{j_t(n)}$. The odd-count transport theorem
identifies $j_t(r)$ with $j_t(n)$, while the shadow congruence identifies
the endpoints modulo $A$. If the displayed congruence held, then
Theorem~\ref{thm:canonical} and $A\ge2^t\ge2$ would force $T^t(r)=1$.

The exact correction identity gives
\[
2^t=A r+d_t(r)\ge2^t r,
\]
so $r\le1$. The zero seed cannot reach one, hence $r=1$. But a positive
return $T^t(1)=1$ has strictly contracting coefficient by the existing
cycle inequality: $A<2^t$. This contradicts the hypothesis.
\end{proof}
\begin{theorem}\label{thm:budget}
Suppose $y\equiv1\pmod{3^K}$, $t>0$, $j_t(n)\le K$, and $T^t(n)=y$.
Then
\[
3^{j_t(n)}<2^t.
\]
In particular, $n$ cannot satisfy $\operatorname{NeverContracts}(n)$.
\end{theorem}
\begin{proof}
Since $3^{j_t(n)}$ divides $3^K$, the target is also one modulo the smaller
ternary scale. If the coefficient did not contract, Theorem~\ref{thm:full}
would give a contradiction. An infinitely noncontracting seed would in
particular have a noncontracting coefficient at time $t$.
\end{proof}

The number of even steps is unrestricted. The odd-step budget matters:
the segment $3\to5\to8\to4$ has two odd steps and coefficient $9/8>1$,
although its endpoint is one modulo three. It is outside the $K=1$ budget.
The positive-length hypothesis also matters; an identity segment need not
contract.

\section{Arbitrarily large unresolved targets}
For any $K$ and any height $B$, choose
\[
y=1+3^{K+1}(B+1).
\]
Then $y>B$, $y\equiv1\pmod3$, and $y\equiv1\pmod{3^K}$.
Theorem~\ref{thm:budget} obstructs every proposed noncontracting predecessor
segment with at most $K$ odd steps. Lean proves this existential statement
for all $K,B$, so the obstruction cannot be removed by excluding just a
fixed finite set of small targets.

This is a restriction on the backward method, not evidence of an escaping
orbit. The constructed targets may all converge; no target is proved to
satisfy $\operatorname{NeverContracts}$. Other arguments could exclude the
residue class directly. The result also does not exclude inverse methods
whose odd-step budget grows with the target, or methods that establish
noncontraction by a different argument.

\section{Formal verification and next obligation}
{\raggedright
Source: \texttt{lean/Collatz/InverseBarrier.lean}. The main declarations are
\texttt{canonical\_endpoint\_lt\_three\_pow},
\texttt{noncontracting\_endpoint\_not\_one\_mod},
\texttt{bounded\_odd\_inverse\_obstruction}, and
\texttt{arbitrarily\_large\_inverse\_obstructions}.
The module build and selected 144-declaration axiom audit pass using only
standard foundations. No custom axiom or \texttt{sorry} is added.
Build with \texttt{lake build Collatz.InverseBarrier} from \texttt{lean/}.\par}

The proposed uniform bounded-odd-step predecessor coverage is therefore
unavailable on this residue class. A full proof still needs an exclusion
of all positive noncontracting seeds and an exclusion of nontrivial cycles.
Neither is supplied here. The Collatz conjecture remains unresolved.
\end{document}
